%% 2-3 ＋ － × ÷ は、何を表しているの？（12ページ、p.38〜49）第7回：見本 ch2-3.tex・ch2-3-mondai.tex を書き直した  % noruby
%% 第1巻（1-4・1-5・1-12・1-14・1-21）で学んだ四則の意味はくり返さず、「量の計算」にしぼる（K7）。  % noruby
%% 第5巻の思い出そうがそのまま成り立つように書く：  % noruby
%%   5-0・5-1「速さ＝進んだ道のり÷かかった時間。変わった量をΔで書く。Δx は x の変化量」  % noruby
%%   5-9「道のり＝速さ×時間。速さが一定なら、v-t グラフの長方形の面積が道のり」  % noruby
%% x・t・v は、この章では「量につけた名前（記号）」として紹介するだけ。文字として計算に使うのは 2-11 から。  % noruby
%% 第11回：ページ制限なし。「なぜ単位をそろえるのか・なぜあと－まえか・時刻と時間・なぜわるのか（2通り）・なぜ m/s×s＝m か」の段と、  % noruby
%%   例題（表から変化量をくらべる・へるときの1分あたり）を足した。  % noruby

%% ---------------- この章だけで使う絵の部品（\bkiid…） ----------------
\makeatletter
% 自転車（横から。右へ進む）  \bkiidbike{(x,y)}{色}   (x,y) は うしろのタイヤの下。長さ 18 mm  % noruby
\newcommand\bkiidbike[2]{\begin{scope}[shift={#1}]
  \draw[bkInk!80, line width=1.1pt] (0,4) circle (4) (14,4) circle (4);
  \draw[#2, line width=1.4pt, line join=round] (0,4) -- (5.5,4) -- (11.5,9.5) -- (4.5,9.5) -- cycle (5.5,4) -- (3.8,11) (11.5,9.5) -- (14,4) (11.5,9.5) -- (12.3,12.5);
  \draw[bkInk!80, line width=1.3pt, line cap=round] (2.4,11.2) -- (5.4,11.2) (11.3,12.6) -- (13.6,12.6);
\end{scope}}
% 人のピクトグラム  \bkiidman{(x,y)}{高さ}{色}   (x,y) は足もと  % noruby
\newcommand\bkiidman[3]{\begin{scope}[shift={#1}, scale=#2/10]
  \fill[#3] (0,8.6) circle (1.35);
  \fill[#3, rounded corners=0.9] (-1.7,0) rectangle (1.7,7);
\end{scope}}
% ストップウォッチ  \bkiidwatch{(x,y)}{半径}   (x,y) は中心  % noruby
\newcommand\bkiidwatch[2]{\begin{scope}[shift={#1}]
  \fill[bkGray!70] (-#2*0.22,#2*0.95) rectangle (#2*0.22,#2*1.3);
  \filldraw[fill=white, draw=bkInk!80, line width=1pt] (0,0) circle (#2);
  \foreach \bkiid@a in {90,0,-90,180} {\draw[bkInk!70, line width=0.6pt] (\bkiid@a:#2*0.82) -- (\bkiid@a:#2*0.62);}
  \draw[bkRed, line width=1pt, line cap=round] (0,0) -- (60:#2*0.7);
  \fill[bkInk] (0,0) circle (#2*0.1);
\end{scope}}
% 旗（コースのスタート・ゴール）  \bkiidflag{(x,y)}{色}   (x,y) は棒の下  % noruby
\newcommand\bkiidflag[2]{\begin{scope}[shift={#1}]
  \draw[bkInk!70, line width=0.8pt] (0,0) -- (0,9);
  \fill[#2] (0,9) -- (4.5,7.6) -- (0,6.2) -- cycle;
\end{scope}}
% 気温の変化（数直線と矢印）  \bkiidtemp{y}{まえ}{あと}{色}{変化のラベル}{左のラベル}   1 ℃＝4.6 mm  % noruby
\newcommand\bkiidtemp[6]{\begin{scope}[shift={(0,#1)}]
  \node[anchor=east, font=\sffamily\bfseries\small, text=bkInk] at (-4,0) {#6};
  \draw[bk axis] (-2,0) -- (96,0);
  \node[anchor=west, bk tick] at (96.5,0) {℃};
  \foreach \bkiid@t in {0,...,20} \draw[bkInk, line width=0.5pt] (\bkiid@t*4.6,-1.2) -- (\bkiid@t*4.6,1.2);
  \foreach \bkiid@t in {0,5,10,15,20} \node[bk tick, font=\sffamily\footnotesize, below=1.6mm] at (\bkiid@t*4.6,0) {\bkiid@t};
  \draw[bk arrow, #4, line width=1.6pt] (#2*4.6,4.5) -- (#3*4.6,4.5);
  \ifdim#3pt>#2pt\node[anchor=west, font=\sffamily\bfseries\small, text=#4!80!black] at (#3*4.6+1.5,4.5) {#5};
  \else\node[anchor=east, font=\sffamily\bfseries\small, text=#4!80!black] at (#3*4.6-1.5,4.5) {#5};\fi
  \node[circle, fill=white, draw=bkGray, line width=1pt, minimum size=3.2mm, inner sep=0pt] at (#2*4.6,0) {};
  \node[circle, fill=#4, minimum size=3.2mm, inner sep=0pt] at (#3*4.6,0) {};
\end{scope}}
% ことばの式を枠で囲んで中央に置く  \bkiidshiki{ことばの式}  % noruby
\newcommand\bkiidshiki[1]{\par\addvspace{3pt}\noindent\hfill
  \tikz[baseline=(s.base)]\node[draw=bkTheme, line width=1pt, fill=white, rounded corners=2mm,
    inner xsep=5mm, inner ysep=1.8mm, font=\sffamily\bfseries] (s) {#1};\hfill\null\par\addvspace{4pt}}
\makeatother

\bkchapter{2}{3}{＋　－　×　÷ は、\\\R{何}{なに}を\R{表}{あらわ}しているの？}{\bklevel{1}}{＋－×÷は\R{何}{なに}を\R{表}{あらわ}しているの？}
\bkmichishirube{1-5,1-12,1-14,2-1,2-2}{2-10,3-5,5-1,5-9}

\begin{goal}
\begin{itemize}[label={\color{bkTheme}●}, leftmargin=1.6\zw, itemsep=0pt, topsep=1pt]
\item \bklevel{1}\enspace \R{同}{おな}じ\R{種類}{しゅ|るい}の\R{量}{りょう}を、\R{単位}{たん|い}をそろえてたしたりひいたりできる。
\item \bklevel{1}\enspace ひき\R{算}{ざん}で\R{変化量}{へん|か|りょう}を\R{求}{もと}め、$\Delta$ を\R{使}{つか}って\R{書}{か}ける。
\item \bklevel{1}\enspace わり\R{算}{ざん}で\R{速}{はや}さを、かけ\R{算}{ざん}で\R{道}{みち}のりを\R{求}{もと}められる。
\end{itemize}
\end{goal}

\begin{hatena}
\R{自転車}{じ|てん|しゃ}が、1\R{秒}{びょう}に\qty{5}{m}ずつ\R{進}{すす}みます。4\R{秒}{びょう}では\R{何}{なん}m\R{進}{すす}むかな？
\begin{bkfigure}
\begin{tikzpicture}[x=1mm, y=1mm]
  \draw[bkGray!60, line width=3pt] (-24,0) -- (100,0);
  \bkiidbike{(-23,1.5)}{bkRoute2}
  \foreach \i in {0,...,4} {
    \pgfmathtruncatemacro\d{5*\i}
    \draw[bkInk, line width=0.8pt] (\i*24,-1.8) -- (\i*24,1.8);
    \node[bk tick, font=\sffamily\footnotesize, below=2mm] at (\i*24,0) {\qty{\d}{m}};
    \fill[bkOrange] (\i*24,6) circle (2);
    \node[bk label, font=\sffamily\footnotesize, above=0.6mm, text=bkOrange!80!black] at (\i*24,8) {\i\R{秒}{びょう}};}
  \foreach \i in {0,...,3}
    \draw[bk arrow, bkGray, line width=1pt] (\i*24+4,6) -- node[above, font=\sffamily\footnotesize, text=bkInk] {\qty{5}{m}} (\i*24+20,6);
\end{tikzpicture}
\end{bkfigure}
\noindent\begin{minipage}[c]{0.5\linewidth}
$5\times4=20$ で、\qty{20}{m}ですね。この\R{進}{すす}み\R{方}{かた}を\R{右}{みぎ}のような\R{図}{ず}にかくと、たて5、よこ4の\R{長方形}{ちょう|ほう|けい}ができます。\R{面積}{めん|せき}も $5\times4=20$。\R{道}{みち}のりと\R{面積}{めん|せき}が、\R{同}{おな}じ20になりました。これは、たまたまかな？
\end{minipage}\hfill
\begin{minipage}[c]{0.47\linewidth}\raggedleft
\begin{tikzpicture}
\begin{axis}[bkplot, width=40mm, height=44mm, xmin=0, xmax=4.8, ymin=0, ymax=6.4, xtick={1,2,3,4}, ytick={1,2,3,4,5},
  xlabel={\R{時間}{じ|かん}（\R{秒}{びょう}）}, ylabel={1\R{秒}{びょう}に\R{進}{すす}む\R{長}{なが}さ（m）}]
  \fill[bkOrange!22] (axis cs:0,0) rectangle (axis cs:4,5);
  \addplot[bkcurve, bkOrange] coordinates {(0,5) (4,5)};
\end{axis}
\end{tikzpicture}
\end{minipage}
\end{hatena}

\clearpage
%% ---------------- p.39 やってみよう・思い出そう ----------------
\begin{yattemiyou}[\R{歩}{ある}く\R{速}{はや}さをくらべよう]
\begin{enumerate}[label={\bkcircled[bkGreen]{\arabic*}}, leftmargin=2.4\zw, itemsep=1pt, topsep=2pt]
\item まきじゃくで、\qty{10}{m}のコースと\qty{15}{m}のコースをつくろう。
\item \R{1人}{ひとり}が\qty{10}{m}、もう\R{1人}{ひとり}が\qty{15}{m}を、ふつうに\R{歩}{ある}こう。かかった\R{時間}{じ|かん}を、\R{友}{とも}だちにストップウォッチではかってもらおう。
\item \R{2人}{ふたり}の\R{歩}{ある}く\R{速}{はや}さは、どちらが\R{速}{はや}いかな？　\R{電卓}{でん|たく}を\R{使}{つか}って\R{考}{かんが}えよう。
\end{enumerate}
\noindent\begin{minipage}[c]{0.6\linewidth}
\begin{tikzpicture}[x=1mm, y=1mm]
  \foreach \y/\L/\c in {15/10/bkRoute2, 0/15/bkRoute4} {
    \pgfmathsetmacro\e{\L*3.4}
    \draw[\c!70, line width=2.2pt] (0,\y) -- (\e,\y);
    \bkiidflag{(0,\y)}{bkGreen}\bkiidflag{(\e,\y)}{bkRed}
    \node[anchor=north, font=\sffamily\bfseries\footnotesize, text=bkInk] at (\e/2,\y-1.5) {\qty{\L}{m}};}
  % 10 m はソラ、15 m はミナ（表と同じ。絵は表2 の2人）  % noruby
  \bkChar{sora}{(-6.5,15)}{12mm}
  \bkChar{mina}{(-6.5,0)}{12mm}
  \node[anchor=east, font=\sffamily\bfseries\footnotesize, text=bkInk] at (-11.5,21) {ソラ};
  \node[anchor=east, font=\sffamily\bfseries\footnotesize, text=bkInk] at (-11.5,6) {ミナ};
  \bkiidwatch{(44.5,20)}{3.6}
\end{tikzpicture}
\end{minipage}\hfill
\begin{minipage}[c]{0.38\linewidth}\centering
\begin{bktable}{colspec={Q[c,m,wd=11mm] Q[c,m,wd=13mm] Q[c,m,wd=17mm]}, row{1}={bg=bkGreen!14, font=\sffamily\bfseries\footnotesize},
  rowsep=2pt, colsep=3pt, cells={font=\small}}
 & \R{道}{みち}のり & かかった\R{時間}{じ|かん} \\
ソラ & \qty{10}{m} & 8\R{秒}{びょう} \\
ミナ & \qty{15}{m} & 10\R{秒}{びょう} \\
\end{bktable}
\par\vspace{2pt}{\footnotesize\sffamily ある\R{班}{はん}の\R{結果}{けっ|か}}
\end{minipage}
\end{yattemiyou}

ソラは、かかった\R{時間}{じ|かん}が\R{短}{みじか}い。でも、\R{歩}{ある}いた\R{道}{みち}のりも\R{短}{みじか}い。ミナは、\R{道}{みち}のりが\R{長}{なが}いけれど、\R{時間}{じ|かん}も\R{長}{なが}い。このままでは、どちらが\R{速}{はや}いかくらべられません。どうすればよいかは、\R{数学}{すう|がく}の\R{道具}{どう|ぐ}③で\R{考}{かんが}えます。

\begin{omoidasou}
\bkrecall{1-5}{ひき\R{算}{ざん}で「のこり」「ちがい」「かわった\R{量}{りょう}（あと－まえ）」を\R{求}{もと}める。}
\bkrecall{1-12}{かけ\R{算}{ざん}は「1つ\R{分}{ぶん}の\R{数}{かず}×いくつ\R{分}{ぶん}」。わり\R{算}{ざん}は\R{同}{おな}じ\R{数}{かず}ずつ\R{分}{わ}けた1つ\R{分}{ぶん}を\R{求}{もと}める（1-14）。}
\bkrecall{1-8}{1\,cm＝10\,mm、1\,m＝100\,cm。\R{重}{おも}さは1\,kg＝1000\,g、かさは1\,L＝1000\,mL（1-9）。}
\bkrecall{1-21}{\R{長方形}{ちょう|ほう|けい}の\R{面積}{めん|せき}＝たて×よこ。\R{単位}{たん|い}は \unit{cm^2} や \unit{m^2}。}
\bkrecall{2-1}{\R{量}{りょう}＝\R{数値}{すう|ち}×\R{単位}{たん|い}。\qty{5}{m} は \qty{1}{m} の5つ\R{分}{ぶん}。}
\bkrecall{2-2}{マイナスは\R{反対}{はん|たい}の\R{向}{む}きを\R{表}{あらわ}す。\R{数直線}{すう|ちょく|せん}の\R{上}{うえ}の\R{動}{うご}きで\R{計算}{けい|さん}できる。}
\end{omoidasou}

\bkfukidashi{mina}{\R{道}{みち}のりも\R{時間}{じ|かん}もちがうと、くらべにくいね。どちらかを\R{同}{おな}じにそろえられたらいいのに。}

%% ---------------- p.40〜43 数学の道具 ----------------
\bksection{doguu}{\R{数学}{すう|がく}の\R{道具}{どう|ぐ}}

\bksubsection{1}{たし\R{算}{ざん}・ひき\R{算}{ざん}は「\R{同}{おな}じ\R{種類}{しゅ|るい}」どうし}

\leavevmode\label{bk:2-3:onaji}\qty{2}{m}のひもと\qty{40}{cm}のひもをつなぐと、\R{何}{なん}cmかな？　\R{数}{かず}だけを\R{見}{み}て $2+40=42$ とすると、42 cm？　42 m？　どちらもおかしいですね。

\begin{bkfigure}
\begin{tikzpicture}[x=1mm, y=1mm]
  \filldraw[fill=bkBlue!22, draw=bkBlue!80!black, line width=0.8pt] (0,0) rectangle (80,7);
  \filldraw[fill=bkOrange!30, draw=bkOrange!80!black, line width=0.8pt] (80,0) rectangle (96,7);
  \node[font=\sffamily\bfseries\small, text=bkInk] at (40,3.5) {\qty{2}{m}（\qty{200}{cm}）};
  \node[font=\sffamily\bfseries\footnotesize, text=bkInk] at (88,3.5) {\qty{40}{cm}};
  \draw[bkInk, line width=0.8pt] (0,-2.5) -- (100,-2.5);
  \foreach \c in {0,50,100,150,200,240} {\draw[bkInk, line width=0.8pt] (\c*0.4,-2.5) -- (\c*0.4,-4.5);
    \node[bk tick, anchor=north, font=\sffamily\footnotesize] at (\c*0.4,-6) {\c};}
  \node[anchor=west, bk tick, font=\sffamily\footnotesize] at (101,-2.5) {cm};
  \draw[decorate, decoration={brace, amplitude=4pt}, line width=0.9pt, bkInk] (0,9) -- (96,9)
    node[midway, above=2mm, font=\sffamily\bfseries\small] {あわせて\qty{240}{cm}};
\end{tikzpicture}
\end{bkfigure}

\R{量}{りょう}は「\R{数値}{すう|ち}×\R{単位}{たん|い}」でした（2-1）。\qty{2}{m}は\qty{1}{m}が2つ、\qty{40}{cm}は\qty{1}{cm}が40こで、1つ\R{分}{ぶん}の\R{大}{おお}きさがちがいます。2と40をそのままたすのは、100\R{円}{えん}\R{玉}{だま}2\R{枚}{まい}と1\R{円}{えん}\R{玉}{だま}40\R{枚}{まい}をあわせて「42\R{枚}{まい}だから42\R{円}{えん}」と\R{言}{い}うようなものです。だから、\R{単位}{たん|い}（1つ\R{分}{ぶん}の\R{大}{おお}きさ）をそろえてからたします。

\textbf{cm にそろえる}\quad 1\,m＝100\,cm（1-8）なので、\qty{2}{m}は100\,cm が2つ\R{分}{ぶん}で、$100\times2=200$ より\qty{200}{cm}です。
\[ \qty{2}{m}+\qty{40}{cm}=\qty{200}{cm}+\qty{40}{cm}=\qty{240}{cm} \]
\textbf{m にそろえても\R{同}{おな}じ}\quad 1\,cm は1\,m を100に\R{分}{わ}けた1つ\R{分}{ぶん}で\qty{0.01}{m}（2-1）。だから\qty{40}{cm}は\qty{0.01}{m}が40こで\qty{0.4}{m}です。
\[ \qty{2}{m}+\qty{40}{cm}=\qty{2}{m}+\qty{0.4}{m}=\qty{2.4}{m} \]
\qty{2.4}{m}と\qty{240}{cm}は、\R{同}{おな}じ\R{長}{なが}さを\R{別}{べつ}の\R{単位}{たん|い}で\R{書}{か}いたものです。どちらにそろえてもかまいません。

\begin{bkcenter}
\begin{bktable}{colspec={Q[l,m,wd=16mm] Q[c,m,wd=48mm] Q[c,m,wd=48mm]}, rowsep=1.5pt, cells={font=\small}}
\R{量}{りょう} & \R{単位}{たん|い}のつながり & \R{例}{れい} \\
\R{長}{なが}さ & 1\,m＝100\,cm、1\,cm＝10\,mm & \qty{3}{m}＝\qty{300}{cm} \\
かさ & 1\,L＝10\,dL＝1000\,mL & \qty{2}{L}＝\qty{2000}{mL} \\
\R{重}{おも}さ & 1\,kg＝1000\,g & \qty{1.5}{kg}＝\qty{1500}{g} \\
\end{bktable}
\bktabcaption{\R{第}{だい}1\R{巻}{かん}（1-8・1-9）で\R{学}{まな}んだ\R{単位}{たん|い}のつながり。}
\end{bkcenter}
\noindent\begin{minipage}[c]{0.62\linewidth}
ひき\R{算}{ざん}も\R{同}{おな}じです。では、「\qty{3}{m}＋2\R{秒}{びょう}」は？　\R{長}{なが}さと\R{時間}{じ|かん}は、\R{単位}{たん|い}をどうそろえても\R{同}{おな}じ\R{種類}{しゅ|るい}の\R{量}{りょう}にならないので、たせません。
\end{minipage}\hfill
\begin{minipage}[c]{0.35\linewidth}\centering
\begin{tikzpicture}[x=1mm, y=1mm]
  \filldraw[fill=bkBlue!22, draw=bkBlue!80!black, line width=0.8pt] (0,0) rectangle (16,6);
  \node[font=\sffamily\bfseries\footnotesize] at (8,3) {\qty{3}{m}};
  \node[font=\sffamily\bfseries\small] at (20.5,3) {＋};
  \bkiidwatch{(29,3)}{4}
  \node[anchor=north, font=\sffamily\bfseries\footnotesize] at (29,-2.5) {2\R{秒}{びょう}};
  \node[font=\sffamily\bfseries\small] at (38,3) {＝};
  \node[font=\sffamily\bfseries\Large, text=bkRed] at (44,3) {？};
\end{tikzpicture}
\end{minipage}

\bksubsection{2}{ひき\R{算}{ざん}で「\R{変化量}{へん|か|りょう}」}

\leavevmode\label{bk:2-3:henka}ある\R{日}{ひ}の\R{気温}{き|おん}は、\R{朝}{あさ}9\R{時}{じ}に\qty{12}{\degreeCelsius}、\R{昼}{ひる}12\R{時}{じ}に\qty{18}{\degreeCelsius}、\R{夜}{よる}9\R{時}{じ}に\qty{10}{\degreeCelsius}でした。\R{何}{なん}℃\R{変}{か}わったかな？

\begin{bkfigure}
\begin{tikzpicture}[x=1mm, y=1mm]
  \bkiidtemp{15}{12}{18}{bkOrange}{$+6$（\R{上}{あ}がった）}{\R{朝}{あさ}→\R{昼}{ひる}}
  \bkiidtemp{0}{18}{10}{bkBlue}{$-8$（\R{下}{さ}がった）}{\R{昼}{ひる}→\R{夜}{よる}}
\end{tikzpicture}
\bkfigcaption{\tikz[baseline=-0.6ex]\node[circle, fill=white, draw=bkGray, line width=1pt, minimum size=2.8mm, inner sep=0pt]{};まえ　\tikz[baseline=-0.6ex]\node[circle, fill=bkGray, minimum size=2.8mm, inner sep=0pt]{};あと。\R{右}{みぎ}へ\R{動}{うご}けばプラス、\R{左}{ひだり}へ\R{動}{うご}けばマイナス。}
\end{bkfigure}

\R{朝}{あさ}→\R{昼}{ひる}は $18-12=6$ で、\qty{6}{\degreeCelsius}\R{上}{あ}がりました。\R{昼}{ひる}→\R{夜}{よる}は $10-18=-8$ です。マイナスは「\R{下}{さ}がった」という\R{反対}{はん|たい}の\R{向}{む}きを\R{表}{あらわ}します（2-2）。

\textbf{なぜ「あと－まえ」なの？}\quad 「まえ」の\R{気温}{き|おん}に、\R{変}{か}わった\R{分}{ぶん}をたすと「あと」の\R{気温}{き|おん}になります。
\[ 12+6=18,\qquad 18+(-8)=10 \]
「まえ＋\R{変}{か}わった\R{分}{ぶん}＝あと」なので、\R{変}{か}わった\R{分}{ぶん}は「あと－まえ」で\R{求}{もと}められます。こうすると、ふえたときはプラス、へったときはマイナスと、\R{向}{む}きも\R{答}{こた}えの\R{中}{なか}に\R{入}{はい}ります。

このように「あと－まえ」で\R{求}{もと}めた、どれだけ\R{変}{か}わったかを\R{表}{あらわ}す\R{量}{りょう}を、\textbf{\R{変化量}{へん|か|りょう}}といいます。\bkidx{へんかりょう}{\R{変化量}{へん|か|りょう}}1-5の「かわった\R{量}{りょう}」に、\R{名前}{な|まえ}がつきました。

\R{変化量}{へん|か|りょう}は、ギリシャ\R{文字}{も|じ}の $\Delta$（デルタ）を\R{使}{つか}って\R{書}{か}きます。\bkidx{でるた}{$\Delta$（デルタ）}$\Delta$ はアルファベットの D にあたる\R{文字}{も|じ}です。「ちがい」を\R{表}{あらわ}す\R{英語}{えい|ご} difference の\R{頭文字}{かしら|も|じ}も d なので、\R{変化量}{へん|か|りょう}のしるしに\R{使}{つか}われています。\R{物理}{ぶつ|り}では、\R{位置}{い|ち}（スタートから\R{何}{なん}mの\R{所}{ところ}か）に $x$、\R{時刻}{じ|こく}に $t$ という\R{名前}{な|まえ}をつけることがよくあります。（\R{文字}{も|じ}の\R{使}{つか}い\R{方}{かた}は、2-11でくわしく\R{学}{まな}びます。）

\begin{bkfigure}
\begin{tikzpicture}[x=1mm, y=1mm]
  \draw[bk axis] (-3,0) -- (124,0);
  \node[anchor=west, bk tick] at (124.5,0) {m};
  \foreach \p in {0,15,...,120} {\draw[bkInk, line width=0.7pt] (\p,-1.4) -- (\p,1.4);
    \node[bk tick, font=\sffamily\footnotesize, below=1.6mm] at (\p,0) {\p};}
  \bkiidbike{(23,2.5)}{bkGray}
  \bkiidbike{(98,2.5)}{bkRoute2}
  \node[anchor=south, font=\sffamily\footnotesize, text=bkGray] at (30,16.5) {まえ（\R{時刻}{じ|こく}10\R{秒}{びょう}）};
  \node[anchor=south, font=\sffamily\footnotesize, text=bkRoute2!80!black] at (105,16.5) {あと（\R{時刻}{じ|こく}25\R{秒}{びょう}）};
  \fill[bkGray] (30,0) circle (1.3);\fill[bkRoute2] (105,0) circle (1.3);
  \draw[decorate, decoration={brace, amplitude=4pt, mirror}, line width=0.9pt, bkRoute2!80!black] (30,-8) -- (105,-8)
    node[midway, below=2mm, font=\sffamily\bfseries\small, text=bkRoute2!80!black] {$\Delta x=\qty{105}{m}-\qty{30}{m}=\qty{75}{m}$};
\end{tikzpicture}
\end{bkfigure}

ここで、「\R{時刻}{じ|こく}」と「\R{時間}{じ|かん}」のちがいに\R{気}{き}をつけましょう。\textbf{\R{時刻}{じ|こく}}は、ストップウォッチがその\R{瞬間}{しゅん|かん}に\R{指}{さ}している\R{数}{かず}（いつ？）。\textbf{\R{時間}{じ|かん}}は、2つの\R{時刻}{じ|こく}のあいだの\R{長}{なが}さ（どれだけのあいだ？）です。

\begin{bkfigure}
\begin{tikzpicture}[x=1mm, y=1mm]
  \draw[bk axis] (-3,0) -- (106,0);
  \node[anchor=west, bk tick] at (106.5,0) {\R{秒}{びょう}};
  \foreach \p in {0,5,...,30} {\draw[bkInk, line width=0.7pt] (\p*3.4,-1.4) -- (\p*3.4,1.4);
    \node[bk tick, font=\sffamily\footnotesize, below=1.6mm] at (\p*3.4,0) {\p};}
  \bkiidwatch{(34,11)}{4.2}
  \bkiidwatch{(85,11)}{4.2}
  \node[anchor=west, font=\sffamily\footnotesize, text=bkGray!80!black] at (39.5,11) {\R{時刻}{じ|こく}10\R{秒}{びょう}};
  \node[anchor=west, font=\sffamily\footnotesize, text=bkRoute2!80!black] at (90.5,11) {\R{時刻}{じ|こく}25\R{秒}{びょう}};
  \fill[bkGray] (34,0) circle (1.3);\fill[bkRoute2] (85,0) circle (1.3);
  \draw[decorate, decoration={brace, amplitude=4pt, mirror}, line width=0.9pt, bkOrange!80!black] (34,-7.5) -- (85,-7.5)
    node[midway, below=2mm, font=\sffamily\bfseries\small, text=bkOrange!80!black] {\R{時間}{じ|かん}（あいだの\R{長}{なが}さ）\quad $25-10=15$ \R{秒}{びょう}};
\end{tikzpicture}
\bkfigcaption{\R{時刻}{じ|こく}は\R{数直線}{すう|ちょく|せん}の\R{上}{うえ}の「\R{点}{てん}」、\R{時間}{じ|かん}は2つの\R{点}{てん}のあいだの「\R{長}{なが}さ」。}
\end{bkfigure}

\R{時刻}{じ|こく}の\R{変化量}{へん|か|りょう}（かかった\R{時間}{じ|かん}）は $\Delta t=\qty{25}{s}-\qty{10}{s}=\qty{15}{s}$（s は\R{秒}{びょう}の\R{記号}{き|ごう}。second の s）、\R{位置}{い|ち}の\R{変化量}{へん|か|りょう}（\R{進}{すす}んだ\R{道}{みち}のり）は $\Delta x=\qty{75}{m}$ です。$\Delta x$ は「$x$ の\R{変化量}{へん|か|りょう}」と\R{読}{よ}みます。$\Delta$ と $x$ のかけ\R{算}{ざん}ではなく、2つで1つの\R{記号}{き|ごう}です。

\bksubsection{3}{わり\R{算}{ざん}で「1あたりの\R{量}{りょう}」}

\leavevmode\label{bk:2-3:wari}\bkyattemiyouref の\R{2人}{ふたり}は、「1\R{秒}{びょう}あたり\R{何}{なん}m\R{進}{すす}んだか」にそろえてくらべます。\qty{10}{m}を8\R{秒}{びょう}で\R{進}{すす}んだなら、\qty{10}{m}を8つに\R{等}{ひと}しく\R{分}{わ}けた1つ\R{分}{ぶん}が、1\R{秒}{びょう}\R{分}{ぶん}の\R{道}{みち}のりです（1-14）。

\begin{bkfigure}
\begin{tikzpicture}[x=1mm, y=1mm]
  \foreach \y/\n/\L/\k/\c/\a in {22/{ソラ\enspace\qty{10}{m}を8\R{秒}{びょう}}/10/8/bkRoute2/{$10\div8=1.25$}, 11/{ミナ\enspace\qty{15}{m}を10\R{秒}{びょう}}/15/10/bkRoute4/{$15\div10=1.5$}, 0/{\R{自転車}{じ|てん|しゃ}\enspace\qty{20}{m}を4\R{秒}{びょう}}/20/4/bkOrange/{$20\div4=5$}} {
    \pgfmathsetmacro\w{\L*3.5}\pgfmathsetmacro\s{\w/\k}
    \fill[\c!15] (0,\y) rectangle (\w,\y+6);
    \fill[\c!70] (0,\y) rectangle (\s,\y+6);
    \foreach \j in {1,...,\k} {\draw[\c!80!black, line width=0.5pt] (\j*\s,\y) -- (\j*\s,\y+6);}
    \draw[\c!80!black, line width=0.8pt] (0,\y) rectangle (\w,\y+6);
    \node[anchor=east, font=\sffamily\bfseries\footnotesize, text=bkInk] at (-2,\y+3) {\n};
    \node[anchor=west, font=\sffamily\small, text=\c!75!black] at (\w+2,\y+3) {\a};}
\end{tikzpicture}
\bkfigcaption{こい\R{色}{いろ}の1つ\R{分}{ぶん}が1\R{秒}{びょう}\R{分}{ぶん}。\R{自転車}{じ|てん|しゃ}は\bkhatenaref の\R{自転車}{じ|てん|しゃ}。}
\end{bkfigure}

$10\div8$ は、\R{電卓}{でん|たく}で\R{計算}{けい|さん}すると1.25です。\R{確}{たし}かめると、$1.25\times8=10$ で、8\R{秒}{びょう}\R{分}{ぶん}あわせるとちゃんと\qty{10}{m}にもどります。ソラは1\R{秒}{びょう}あたり\qty{1.25}{m}、ミナは\qty{1.5}{m}。ミナのほうが1\R{秒}{びょう}に\R{長}{なが}く\R{進}{すす}むので、\R{速}{はや}いとわかります。どれも、\R{道}{みち}のりを\R{時間}{じ|かん}でわっています。この「1\R{秒}{びょう}あたりに\R{進}{すす}む\R{道}{みち}のり」を、\textbf{\R{速}{はや}さ}といいます。\bkidx{はやさ}{\R{速}{はや}さ}
\bkiidshiki{\R{速}{はや}さ＝\R{進}{すす}んだ\R{道}{みち}のり÷かかった\R{時間}{じ|かん}}
\textbf{もう1つのそろえ\R{方}{かた}}\quad \R{反対}{はん|たい}に、「1\,m あたり\R{何秒}{なん|びょう}かかったか」にそろえても、くらべられます。ソラは $8\div10=0.8$ で1\,m に0.8\R{秒}{びょう}、ミナは $10\div15=0.66\cdots$ で1\,m に\R{約}{やく}0.67\R{秒}{びょう}。1\,m に\R{短}{みじか}い\R{時間}{じ|かん}しかかからないミナのほうが\R{速}{はや}く、\R{答}{こた}えは\R{同}{おな}じです。でもこの\R{数}{かず}は、\R{小}{ちい}さいほど\R{速}{はや}いので、まちがえやすいですね。そこで、\R{数}{かず}が\R{大}{おお}きいほど\R{速}{はや}い「1\R{秒}{びょう}あたりの\R{道}{みち}のり」を、\R{速}{はや}さとよぶことにしています。
1\R{秒}{びょう}あたり\qty{1.5}{m}の\R{速}{はや}さを、\qty{1.5}{m/s}（1.5メートル\R{毎秒}{まい|びょう}）と\R{書}{か}きます。s は\R{秒}{びょう}、/ は÷のしるしです（1-14のコラム）。\R{単位}{たん|い}も\R{数}{かず}といっしょにわり\R{算}{ざん}できます。
\[ \qty{15}{m}\div\qty{10}{s}=\qty{1.5}{m/s} \]
②の\R{自転車}{じ|てん|しゃ}の\R{速}{はや}さは、$\Delta x\div\Delta t=\qty{75}{m}\div\qty{15}{s}=\qty{5}{m/s}$ です。とちゅうで\R{速}{はや}さが\R{変}{か}わっても、\R{道}{みち}のり÷\R{時間}{じ|かん}で\R{出}{だ}した\R{速}{はや}さは\R{全体}{ぜん|たい}をならした\R{速}{はや}さで、\textbf{\R{平均}{へい|きん}の\R{速}{はや}さ}といいます（\R{平均}{へい|きん}は2-6でくわしく\R{学}{まな}びます）。

\noindent\begin{minipage}[c]{0.5\linewidth}
\R{速}{はや}さのほかにも、「1あたり」にそろえるとくらべられる\R{量}{りょう}があります。\R{遊}{あそ}び\R{場}{ば}Aは\qty{6}{m^2}に12\R{人}{にん}、Bは\qty{8}{m^2}に20\R{人}{にん}。どちらがこんでいるかな？\par
\R{広}{ひろ}さがちがうので、1\,\unit{m^2}あたりの\R{人数}{にん|ずう}にそろえます。\R{人数}{にん|ずう}を\R{広}{ひろ}さでわります。\par
A：$12\div6=2$\quad B：$20\div8=2.5$
\end{minipage}\hfill
\begin{minipage}[c]{0.47\linewidth}\centering
\begin{tikzpicture}[x=1mm, y=1mm]
  \foreach \i in {0,1,2} {\foreach \j in {0,1} {
    \filldraw[fill=bkRoute4!10, draw=bkRoute4!70!black, line width=0.5pt] (\i*8.5,\j*8.5) rectangle ++(8.5,8.5);
    \foreach \dx in {2.8,5.7} {\bkiidman{(\i*8.5+\dx,\j*8.5+1.4)}{5.4}{bkRoute4!80!black}}}}
  \node[anchor=south, font=\sffamily\bfseries\footnotesize] at (12.75,18) {A};
  \begin{scope}[shift={(30,0)}]
  \foreach \i/\j/\n in {0/0/3,1/0/2,2/0/3,3/0/2,0/1/2,1/1/3,2/1/2,3/1/3} {
    \filldraw[fill=bkRoute2!10, draw=bkRoute2!70!black, line width=0.5pt] (\i*8.5,\j*8.5) rectangle ++(8.5,8.5);
    \ifnum\n=3 {\foreach \dx in {2,4.25,6.5} {\bkiidman{(\i*8.5+\dx,\j*8.5+1.4)}{5.4}{bkRoute2!80!black}}}
    \else {\foreach \dx in {2.8,5.7} {\bkiidman{(\i*8.5+\dx,\j*8.5+1.4)}{5.4}{bkRoute2!80!black}}}\fi}
  \node[anchor=south, font=\sffamily\bfseries\footnotesize] at (17,18) {B};
  \end{scope}
  \node[anchor=north, font=\sffamily\footnotesize] at (32,-1.5) {1マスが\qty{1}{m^2}};
\end{tikzpicture}
\end{minipage}
\par\vspace{2pt}
1\,\unit{m^2}あたりの\R{人数}{にん|ずう}が\R{多}{おお}いBのほうが、こんでいます。この「1\,\unit{m^2}あたりの\R{人数}{にん|ずう}」を\textbf{こみぐあい}といいます（\R{単位}{たん|い}は\R{人}{にん}/\unit{m^2}）。\bkidx{こみぐあい}{こみぐあい}\R{速}{はや}さもこみぐあいも、わり\R{算}{ざん}でつくる「1あたりの\R{量}{りょう}」です。

\bksubsection{4}{かけ\R{算}{ざん}で「\R{新}{あたら}しい\R{量}{りょう}」}

\leavevmode\label{bk:2-3:kake}\bkhatenaref の\R{自転車}{じ|てん|しゃ}は、\R{速}{はや}さ\qty{5}{m/s}です。4\R{秒}{びょう}では、1\R{秒}{びょう}\R{分}{ぶん}の\qty{5}{m}が4つ\R{分}{ぶん}。いくつ\R{分}{ぶん}かを「\R{何倍}{なん|ばい}」ともいうので、\qty{5}{m}の4\R{倍}{ばい}です。
\[ \qty{5}{m/s}\times\qty{4}{s}=\qty{20}{m} \]
③の\R{反対}{はん|たい}に、「1あたり」から\R{全体}{ぜん|たい}をつくる\R{計算}{けい|さん}です。

\textbf{\R{単位}{たん|い}はどうなる？}\quad m/s は「m÷s」でした。それに s をかけると、\mbox{「m÷s×s」}です。\R{数}{かず}で\R{考}{かんが}えると、$20\div4\times4=5\times4=20$ のように、\R{同}{おな}じ\R{数}{かず}でわってかけると\R{元}{もと}にもどります。\R{単位}{たん|い}も\R{同}{おな}じで、÷s と×s で\R{元}{もと}にもどり、m だけがのこります。だから\R{答}{こた}えの\R{単位}{たん|い}は、\R{道}{みち}のりの m になります。
\[ \unit{m/s}\times\unit{s}=\unit{m} \]
\bkiidshiki{\R{道}{みち}のり＝\R{速}{はや}さ×\R{時間}{じ|かん}}

\noindent\begin{minipage}[c]{0.5\linewidth}
\begin{tikzpicture}
\begin{axis}[bkplot, width=56mm, height=44mm, xmin=0, xmax=4.9, ymin=0, ymax=6.6, xtick={1,2,3,4}, ytick={1,2,3,4,5},
  xlabel={\R{時間}{じ|かん}（s）}, ylabel={\R{速}{はや}さ（m/s）}]
  \fill[bkBlue!16] (axis cs:0,0) rectangle (axis cs:4,5);
  \fill[bkBlue!40] (axis cs:0,0) rectangle (axis cs:1,5);
  \addplot[bkcurve] coordinates {(0,5) (4,5)};
  \foreach \k in {0,1,2,3} {\edef\tmp{\noexpand\node[font=\noexpand\sffamily\noexpand\bfseries\noexpand\footnotesize, text=bkBlue!80!black, fill=white, inner sep=1pt] at (axis cs:\k+0.5,2.5) {\noexpand\qty{5}{m}};}\tmp}
\end{axis}
\end{tikzpicture}
\end{minipage}\hfill
\begin{minipage}[c]{0.47\linewidth}
\R{横}{よこ}に\R{時間}{じ|かん}、たてに\R{速}{はや}さをとります。こい\R{色}{いろ}の\R{細}{ほそ}い\R{長方形}{ちょう|ほう|けい}（たて5、よこ1）の\R{面積}{めん|せき} $5\times1=5$ は、1\R{秒}{びょう}\R{間}{かん}に\R{進}{すす}む\qty{5}{m}。4つならべた\R{長方形}{ちょう|ほう|けい}の\R{面積}{めん|せき} $5\times4=20$ は、4\R{秒}{びょう}\R{間}{かん}の\R{道}{みち}のり\qty{20}{m}です。
\end{minipage}
\par\vspace{3pt}
\R{速}{はや}さ×\R{時間}{じ|かん}も、たて×よこも、\R{同}{おな}じかけ\R{算}{ざん}です。だから、\R{速}{はや}さが\R{一定}{いっ|てい}なら、\R{速}{はや}さと\R{時間}{じ|かん}のグラフの\R{長方形}{ちょう|ほう|けい}の\R{面積}{めん|せき}が、\R{進}{すす}んだ\R{道}{みち}のりになります。\bkhatenaref の\R{答}{こた}えは「たまたまではない」です。\R{速}{はや}さを $v$、\R{時刻}{じ|こく}を $t$ で\R{表}{あらわ}すことが\R{多}{おお}いので、このグラフを\textbf{$v$-$t$ グラフ}ともいいます。\bkidx{ぶいてぃーぐらふ}{$v$-$t$ グラフ}

\begin{taisetsu}
＋－\enspace \R{同}{おな}じ\R{種類}{しゅ|るい}の\R{量}{りょう}どうし。\R{単位}{たん|い}をそろえてから\R{計算}{けい|さん}する。\\
－\enspace \R{変化量}{へん|か|りょう}＝あと－まえ（$\Delta$ で\R{書}{か}く。マイナスにもなる）\\
÷\enspace \R{速}{はや}さ＝\R{進}{すす}んだ\R{道}{みち}のり÷かかった\R{時間}{じ|かん}（1あたりの\R{量}{りょう}）\\
×\enspace \R{道}{みち}のり＝\R{速}{はや}さ×\R{時間}{じ|かん}。\R{速}{はや}さが\R{一定}{いっ|てい}なら、\\
\hspace*{1.5\zw}\R{速}{はや}さと\R{時間}{じ|かん}のグラフの\R{長方形}{ちょう|ほう|けい}の\R{面積}{めん|せき}が\R{道}{みち}のり
\end{taisetsu}

\noindent\begin{minipage}[c]{0.5\linewidth}
\begin{kotoba}
\small
\bkword{\R{変化量}{へん|か|りょう}}{あと－まえ。$\Delta$ で\R{書}{か}く。}
\bkword{\R{速}{はや}さ}{1\R{秒}{びょう}あたりに\R{進}{すす}む\R{道}{みち}のり。m/s（メートル\R{毎秒}{まい|びょう}）。}
\bkword{こみぐあい}{1\,\unit{m^2}あたりの\R{人数}{にん|ずう}。}
\bkword{$v$-$t$ グラフ}{\R{速}{はや}さと\R{時間}{じ|かん}のグラフ。}
\end{kotoba}
\end{minipage}\hfill
\begin{minipage}[c]{0.48\linewidth}\raggedleft
\bkkigocard{$\Delta$}{デルタ}{\R{変化量}{へん|か|りょう}。「あと－まえ」で\R{求}{もと}める。$\Delta x$ は $x$ の\R{変化量}{へん|か|りょう}。}{delta}{2-3}
\end{minipage}

\bkfukidashi{delta}{$\Delta$ は、ぼくの\R{名前}{な|まえ}と\R{同}{おな}じ「デルタ」。\R{第}{だい}5\R{巻}{かん}の\R{微分}{び|ぶん}でも、\R{毎}{まい}ページのように\R{出}{で}てくるしるしだよ。}

%% ---------------- 例題 ----------------
\bksection{reidai}{\R{例題}{れい|だい}でたしかめよう}

\begin{reidai}[たせる？　たせない？]
\R{計算}{けい|さん}できるかな？\quad（1）\nobreak $\qty{3}{m}+\qty{40}{cm}$\quad（2）\nobreak $\qty{1}{L}-\qty{300}{mL}$\quad（3）\nobreak \mbox{\qty{5}{m}＋3\R{秒}{びょう}}\par
\kangaekata （1）は\R{長}{なが}さどうし。1\,m＝100\,cm を\R{使}{つか}って cm にそろえる。（2）はかさどうし。1\,L＝1000\,mL を\R{使}{つか}って mL にそろえる。（3）は\R{長}{なが}さと\R{時間}{じ|かん}。\par
\shiki （1）\nobreak $\qty{3}{m}=\qty{300}{cm}$ なので、$\qty{300}{cm}+\qty{40}{cm}=\qty{340}{cm}$\par
\hspace*{2.6\zw}（2）\nobreak $\qty{1}{L}=\qty{1000}{mL}$ なので、$\qty{1000}{mL}-\qty{300}{mL}=\qty{700}{mL}$\par
\kotae （1）\nobreak \qty{340}{cm}\hspace{1.5\zw}（2）\nobreak \qty{700}{mL}\hspace{1.5\zw}（3）\nobreak \R{種類}{しゅ|るい}がちがうので\R{計算}{けい|さん}できない
\end{reidai}

\begin{jibun}
\R{計算}{けい|さん}できる？\enspace（1）\nobreak $\qty{1}{kg}+\qty{250}{g}$\enspace（2）\nobreak $\qty{2}{m}-\qty{60}{cm}$\enspace（3）\nobreak 10\R{分}{ぷん}＋\qty{10}{m}
\bkjans{（1）\nobreak \qty{1250}{g}\quad（2）\nobreak \qty{140}{cm}\quad（3）\nobreak \R{時間}{じ|かん}と\R{長}{なが}さなので\R{計算}{けい|さん}できない}
\end{jibun}

\begin{reidai}[\R{変化量}{へん|か|りょう}を\R{求}{もと}める]
\R{気温}{き|おん}が、\R{朝}{あさ}\qty{-2}{\degreeCelsius}→\R{昼}{ひる}\qty{7}{\degreeCelsius}→\R{夜}{よる}\qty{1}{\degreeCelsius}と\R{変}{か}わりました。（1）\nobreak \R{朝}{あさ}→\R{昼}{ひる}、（2）\nobreak \R{昼}{ひる}→\R{夜}{よる}の\R{変化量}{へん|か|りょう}は？\par
\kangaekata あと－まえ。（1）は\R{数直線}{すう|ちょく|せん}で $-2$ から7まで\R{右}{みぎ}へ9\R{進}{すす}む（2-2）。\par
\shiki （1）\nobreak $7-(-2)=9$\hspace{2\zw}（2）\nobreak $1-7=-6$\par
\kotae （1）\nobreak $+\qty{9}{\degreeCelsius}$（\R{上}{あ}がった）\hspace{1.5\zw}（2）\nobreak $\qty{-6}{\degreeCelsius}$（\R{下}{さ}がった）
\end{reidai}

\begin{jibun}
\R{気温}{き|おん}が、\qty{-5}{\degreeCelsius}→\qty{3}{\degreeCelsius}→\qty{-1}{\degreeCelsius}と\R{変}{か}わりました。2つの\R{変化量}{へん|か|りょう}を\R{求}{もと}めよう。
\bkjans{$3-(-5)=8$\quad $+\qty{8}{\degreeCelsius}$（\R{上}{あ}がった）、$-1-3=-4$\quad $\qty{-4}{\degreeCelsius}$（\R{下}{さ}がった）}
\end{jibun}

\begin{reidai}[\R{表}{ひょう}から\R{変化量}{へん|か|りょう}をくらべる]
\R{表}{ひょう}は、ある\R{日}{ひ}の\R{気温}{き|おん}です。\R{気温}{き|おん}がいちばん\R{大}{おお}きく\R{上}{あ}がったのは、どの2\R{時間}{じ|かん}でしょう。\par
\begin{bkcenter}
\begin{bktable}{colspec={Q[l,m,wd=24mm] *{4}{Q[c,m,wd=15mm]}}, rowsep=1.5pt}
\R{時刻}{じ|こく} & 8\R{時}{じ} & 10\R{時}{じ} & 12\R{時}{じ} & 14\R{時}{じ} \\
\R{気温}{き|おん}（℃） & 4 & 9 & 12 & 11 \\
\end{bktable}
\end{bkcenter}
\kangaekata となりどうしの\R{時刻}{じ|こく}で、1つずつ「あと－まえ」を\R{計算}{けい|さん}して、\R{表}{ひょう}に\R{書}{か}きこむ。プラスで、いちばん\R{大}{おお}きいものをえらぶ。\par
\shiki 8\R{時}{じ}→10\R{時}{じ}：$9-4=5$\quad 10\R{時}{じ}→12\R{時}{じ}：$12-9=3$\par
\hspace*{2.6\zw}12\R{時}{じ}→14\R{時}{じ}：$11-12=-1$\par
\hspace*{2.6\zw}いちばん\R{大}{おお}きいのは $+5$\par
\kotae 8\R{時}{じ}から10\R{時}{じ}（$+\qty{5}{\degreeCelsius}$）
\end{reidai}

\begin{jibun}
\R{水}{みず}そうの\R{水}{みず}の\R{量}{りょう}が、0\R{分}{ふん}に\qty{10}{L}、2\R{分}{ふん}に\qty{16}{L}、4\R{分}{ふん}に\qty{25}{L}、6\R{分}{ふん}に\qty{31}{L}でした。いちばん\R{多}{おお}くふえたのは、どの2\R{分}{ふん}\R{間}{かん}でしょう。
\bkjans{0→2\R{分}{ふん}は $16-10=6$、2→4\R{分}{ふん}は $25-16=9$、4→6\R{分}{ふん}は $31-25=6$。2\R{分}{ふん}から4\R{分}{ふん}（$+\qty{9}{L}$）}
\end{jibun}

\begin{reidai}[$\Delta$ を\R{使}{つか}って\R{速}{はや}さを\R{求}{もと}める]
\R{電車}{でん|しゃ}が、\R{時刻}{じ|こく}20\R{秒}{びょう}にスタートから\qty{150}{m}、\R{時刻}{じ|こく}50\R{秒}{びょう}に\qty{750}{m}の\R{位置}{い|ち}にいました。$\Delta t$ と $\Delta x$ を\R{求}{もと}めて、\R{速}{はや}さを\R{出}{だ}そう。\par
\kangaekata $\Delta t$ も $\Delta x$ も「あと－まえ」。\R{速}{はや}さは $\Delta x\div\Delta t$。\par
\shiki $\Delta t=\qty{50}{s}-\qty{20}{s}=\qty{30}{s}$、$\Delta x=\qty{750}{m}-\qty{150}{m}=\qty{600}{m}$\par
\hspace*{2.6\zw}$\Delta x\div\Delta t=\qty{600}{m}\div\qty{30}{s}=\qty{20}{m/s}$\kotaenext\R{速}{はや}さ\qty{20}{m/s}
\end{reidai}

\begin{jibun}
\R{水}{みず}そうに\R{水}{みず}を\R{入}{い}れています。\R{時刻}{じ|こく}5\R{分}{ふん}に\qty{25}{L}、\R{時刻}{じ|こく}9\R{分}{ふん}に\qty{61}{L}でした。\R{時刻}{じ|こく}の\R{変化量}{へん|か|りょう}と\R{水}{みず}の\R{量}{りょう}の\R{変化量}{へん|か|りょう}を\R{求}{もと}め、1\R{分}{ぷん}あたり\R{何}{なん}Lふえたかを\R{出}{だ}そう。
\bkjans{\R{時刻}{じ|こく}の\R{変化量}{へん|か|りょう} $9-5=4$（\R{分}{ふん}）、\R{水}{みず}の\R{量}{りょう}の\R{変化量}{へん|か|りょう} $61-25=36$（L）、$36\div4=9$\quad 1\R{分}{ぷん}あたり\qty{9}{L}}
\end{jibun}

\begin{reidai}[へっていくときの1あたり]
\R{雨水}{あま|みず}をためたタンクの\R{水}{みず}が、\R{時刻}{じ|こく}0\R{分}{ふん}に\qty{50}{L}、\R{時刻}{じ|こく}5\R{分}{ふん}に\qty{35}{L}でした。\R{水}{みず}の\R{量}{りょう}の\R{変化量}{へん|か|りょう}と、1\R{分}{ぷん}あたり\R{何}{なん}Lへったかを\R{求}{もと}めよう。\par
\kangaekata \R{変化量}{へん|か|りょう}は「あと－まえ」。へったので、マイナスになるはず。「1\R{分}{ぷん}あたり\R{何}{なん}Lへったか」は、へった\R{量}{りょう}（マイナスをとった\R{大}{おお}きさ）を\R{時間}{じ|かん}でわる。\par
\shiki \R{水}{みず}の\R{量}{りょう}の\R{変化量}{へん|か|りょう}：$35-50=-15$（L）\quad\R{時間}{じ|かん}：$5-0=5$（\R{分}{ふん}）\par
\hspace*{2.6\zw}\qty{15}{L}へったので、$15\div5=3$\par
\kotae \R{変化量}{へん|か|りょう}は $\qty{-15}{L}$（\qty{15}{L}へった）、1\R{分}{ぷん}あたり\qty{3}{L}へった
\end{reidai}

\begin{jibun}
ろうそくの\R{長}{なが}さが、\R{時刻}{じ|こく}0\R{分}{ふん}に\qty{12}{cm}、\R{時刻}{じ|こく}20\R{分}{ぷん}に\qty{8}{cm}でした。\R{長}{なが}さの\R{変化量}{へん|か|りょう}と、1\R{分}{ぷん}あたり\R{何}{なん}cmへったかを\R{求}{もと}めよう。
\bkjans{$8-12=-4$\quad $\qty{-4}{cm}$（\qty{4}{cm}へった）。$4\div20=0.2$\quad 1\R{分}{ぷん}あたり\qty{0.2}{cm}へった}
\end{jibun}

\begin{reidai}[\R{速}{はや}さ×\R{時間}{じ|かん}と、グラフの\R{面積}{めん|せき}]
\R{自転車}{じ|てん|しゃ}が\qty{3}{m/s}で2\R{秒}{びょう}、\qty{2}{m/s}で4\R{秒}{びょう}\R{走}{はし}りました。\R{何}{なん}m\R{進}{すす}んだかな？\par
\kangaekata グラフの2つの\R{長方形}{ちょう|ほう|けい}の\R{面積}{めん|せき}（\R{速}{はや}さ×\R{時間}{じ|かん}）をたす。\par
\noindent\begin{minipage}[c]{0.5\linewidth}
\shiki $3\times2=6$、\enspace$2\times4=8$\par
\hspace*{2.6\zw}$6+8=14$\kotaenext\qty{14}{m}
\end{minipage}\hfill
\begin{minipage}[c]{0.48\linewidth}\raggedleft
\begin{tikzpicture}
\begin{axis}[bkplot, width=50mm, height=30mm, xmin=0, xmax=6.6, ymin=0, ymax=3.6, xtick={2,6}, ytick={2,3},
  xlabel={\R{時間}{じ|かん}（s）}, ylabel={\R{速}{はや}さ（m/s）}]
  \fill[bkBlue!20] (axis cs:0,0) rectangle (axis cs:2,3);
  \fill[bkOrange!28] (axis cs:2,0) rectangle (axis cs:6,2);
  \addplot[bkcurve] coordinates {(0,3) (2,3)};
  \addplot[bkcurve] coordinates {(2,2) (6,2)};
  \node[font=\sffamily\bfseries\footnotesize, text=bkBlue!80!black] at (axis cs:1,1.5) {\qty{6}{m}};
  \node[font=\sffamily\bfseries\footnotesize, text=bkOrange!80!black] at (axis cs:4,1) {\qty{8}{m}};
\end{axis}
\end{tikzpicture}
\end{minipage}
\end{reidai}

\begin{jibun}
エアコンが、1\R{時間}{じ|かん}あたり\qty{1.2}{kWh}で2\R{時間}{じ|かん}、そのあと\qty{0.4}{kWh}で3\R{時間}{じ|かん}\R{使}{つか}いました。\R{全部}{ぜん|ぶ}で\R{何}{なん}kWhでしょう。（kWh は\R{電気}{でん|き}の\R{量}{りょう}の\R{単位}{たん|い}）
\bkjans{$1.2\times2=2.4$、$0.4\times3=1.2$、$2.4+1.2=3.6$\quad\qty{3.6}{kWh}}
\end{jibun}

%% ---------------- まちがえやすいところ ----------------
\bksection{tsumazuki}{まちがえやすいところをチェック}

\begin{tsumazuki}{\R{単位}{たん|い}をそろえずにたす（\qty{2}{m}のひもと\qty{40}{cm}のひもをつなぐ）}
\bkng{$\qty{2}{m}+\qty{40}{cm}=42$\quad（42 cm？　42 m？）}
\bkok{$\qty{2}{m}+\qty{40}{cm}=\qty{200}{cm}+\qty{40}{cm}=\qty{240}{cm}$}
1つ\R{分}{ぶん}（\R{単位}{たん|い}）がちがう\R{数}{かず}は、そのままたせません。\R{答}{こた}えにいつも\R{単位}{たん|い}を\R{書}{か}くと、まちがいに\R{気}{き}づけます。
\end{tsumazuki}

\begin{tsumazuki}{ひき\R{算}{ざん}の\R{順番}{じゅん|ばん}がぎゃく（\qty{80}{\degreeCelsius}のお\R{湯}{ゆ}が\qty{65}{\degreeCelsius}に）}
\bkng{まえ－あと\quad $80-65=15$\quad（\qty{15}{\degreeCelsius}\R{上}{あ}がった？）}
\bkok{あと－まえ\quad $65-80=-15$\quad（\qty{15}{\degreeCelsius}\R{下}{さ}がった）}
ぎゃくにすると、プラスとマイナスが\R{反対}{はん|たい}になり、\R{冷}{さ}めたのに「\R{上}{あ}がった」ことになります。
\end{tsumazuki}

\begin{tsumazuki}{わる\R{順番}{じゅん|ばん}がぎゃく（\qty{20}{m}を4\R{秒}{びょう}で\R{進}{すす}んだ\R{速}{はや}さ）}
\bkng{$4\div20=0.2$\quad（\qty{0.2}{m/s}？）}
\bkok{$20\div4=5$\quad（\qty{5}{m/s}）}
\R{速}{はや}さは「1\R{秒}{びょう}あたり」の\R{道}{みち}のり。「○○あたり」の○○（ここでは\R{秒}{びょう}）でわります。
\end{tsumazuki}

%% ---------------- 物理・工学ではこう使う ----------------
\begin{butsuri}
\noindent\begin{minipage}[c]{0.3\linewidth}\centering
\begin{tikzpicture}[x=1mm, y=1mm]
  \fill[bkInk!85, rounded corners=2.5mm] (0,0) rectangle (26,44);
  \fill[bkGreen!12] (2,5) rectangle (24,39);
  \draw[white, line width=2.2pt] (2,11) -- (24,33);
  \draw[white, line width=1.6pt] (13,5) -- (13,39);
  \fill[bkBlue] (7,16) circle (1.5);
  \fill[bkRed] (17,26) circle (1.5);
  \draw[bkRed, line width=0.9pt, -{Stealth[length=3pt]}] (8.4,17.4) -- (15.6,24.6);
  \node[anchor=north west, font=\sffamily\bfseries\scriptsize, text=bkInk, fill=white, inner sep=1pt] at (3,38) {\qty{3.0}{m/s}};
  \fill[bkGray!60] (10,1.6) rectangle (16,3);
\end{tikzpicture}
\end{minipage}\hfill
\begin{minipage}[c]{0.67\linewidth}
\textbf{スマートフォンの\R{速}{はや}さの\R{表示}{ひょう|じ}}\par
スマートフォンやランニング\R{用}{よう}の\R{時計}{とけい}は、\R{人工衛星}{じん|こう|えい|せい}からの\R{電波}{でん|ぱ}で、\R{自分}{じ|ぶん}の\R{位置}{い|ち}を1\R{秒}{びょう}ごとにはかれます（GPS）。1\R{秒}{びょう}で\R{位置}{い|ち}が\qty{3}{m}\R{変}{か}わったら、$\Delta x\div\Delta t=3\div1=3$ で、\R{速}{はや}さは\qty{3}{m/s}と\R{計算}{けい|さん}できます。
\end{minipage}

\par\vspace{6pt}
\noindent\begin{minipage}[c]{0.3\linewidth}\centering
\begin{tikzpicture}[x=1mm, y=1mm]
  \filldraw[fill=white, draw=bkGray, line width=1pt, rounded corners=2mm] (0,0) rectangle (32,26);
  \fill[bkInk!85, rounded corners=1mm] (3,12) rectangle (29,22);
  \node[font=\sffamily\bfseries\small, text=bkGreen!30] at (16,17) {1234.5};
  \node[anchor=north, font=\sffamily\bfseries\scriptsize, text=bkInk] at (16,10) {kWh};
  \fill[bkGray!50] (12,-4) rectangle (20,0);
\end{tikzpicture}
\end{minipage}\hfill
\begin{minipage}[c]{0.67\linewidth}
\textbf{\R{家}{いえ}の\R{電気}{でん|き}のメーター}\par
\R{使}{つか}った\R{電気}{でん|き}の\R{量}{りょう}（\R{電力量}{でん|りょく|りょう}）は、kWh（キロワット\R{時}{じ}）という\R{単位}{たん|い}ではかります。1\R{時間}{じ|かん}あたり\qty{0.6}{kWh}を\R{使}{つか}うエアコンを5\R{時間}{じ|かん}つけると、$0.6\times5=3$ で\qty{3}{kWh}。「1あたり×いくつ\R{分}{ぶん}」で、\R{電気}{でん|き}の\R{量}{りょう}ができます。
\end{minipage}
\end{butsuri}

\begin{column}{＋－×÷は、どこへつながる？}
この\R{章}{しょう}の「ひき\R{算}{ざん}・わり\R{算}{ざん}」と「かけ\R{算}{ざん}」は、\R{工学}{こう|がく}のいちばん\R{大事}{だい|じ}な\R{道具}{どう|ぐ}に\R{育}{そだ}ちます。
\begin{bkfigure}
\begin{tikzpicture}[x=1mm, y=1mm]
  \node[bk node, minimum width=28mm] (hiki) at (0,19) {$\Delta x\div\Delta t$\\{\footnotesize この\R{章}{しょう}}};
  \node[bk node, minimum width=28mm] (kake) at (0,0) {\R{速}{はや}さ×\R{時間}{じ|かん}\\{\footnotesize この\R{章}{しょう}}};
  \node[bk node, draw=bkPurple, minimum width=26mm] (kata) at (35,19) {グラフの\R{傾}{かたむ}き\\{\footnotesize 3-5}};
  \node[bk node, draw=bkPurple, minimum width=26mm] (men) at (35,0) {グラフの\R{面積}{めん|せき}\\{\footnotesize 5-9}};
  \node[bk node, draw=bkPurple, minimum width=18mm] (bibun) at (66,19) {\textbf{\R{微分}{び|ぶん}}\\{\footnotesize \R{第}{だい}5\R{巻}{かん}}};
  \node[bk node, draw=bkPurple, minimum width=18mm] (sekibun) at (66,0) {\textbf{\R{積分}{せき|ぶん}}\\{\footnotesize \R{第}{だい}5\R{巻}{かん}}};
  \node[bk node, draw=bkPurple, fill=bkPurple!10] (goal) at (100,9.5) {\textbf{\R{微分方程式}{び|ぶん|ほう|てい|しき}}\\{\footnotesize \R{第}{だい}6\R{巻}{かん}}};
  \draw[bk link] (hiki) -- (kata);
  \draw[bk link] (kake) -- (men);
  \draw[bk link] (kata) -- (bibun);
  \draw[bk link] (men) -- (sekibun);
  \draw[bk link] (bibun) -- (goal);
  \draw[bk link] (sekibun) -- (goal);
\end{tikzpicture}
\end{bkfigure}
\R{微分方程式}{び|ぶん|ほう|てい|しき}を\R{使}{つか}うと、ボールの\R{動}{うご}きやお\R{湯}{ゆ}の\R{冷}{さ}め\R{方}{かた}を、\R{式}{しき}で\R{表}{あらわ}して\R{計算}{けい|さん}できます。その\R{始}{はじ}まりが、この\R{章}{しょう}の「あと－まえ」と「\R{長方形}{ちょう|ほう|けい}の\R{面積}{めん|せき}」です。
\end{column}

\clearpage
%% ---------------- 章末問題 ----------------
\bkendproblems
\bkproblemgroup{\R{基本}{き|ほん}}{1}

\bkproblem \R{次}{つぎ}の\R{計算}{けい|さん}をしましょう。\R{答}{こた}えには\R{単位}{たん|い}をつけましょう。\par
（1）\nobreak $\qty{3}{m}+\qty{25}{cm}$\hspace{1.5\zw}（2）\nobreak $\qty{2}{L}-\qty{500}{mL}$\hspace{1.5\zw}（3）\nobreak $\qty{6}{m/s}\times\qty{5}{s}$\hspace{1.5\zw}（4）\nobreak $\qty{36}{m}\div\qty{9}{s}$
\bkpans{（1）\nobreak \qty{325}{cm}\quad（2）\nobreak \qty{1500}{mL}\quad（3）\nobreak \qty{30}{m}\quad（4）\nobreak \qty{4}{m/s}}

\bkproblem \R{計算}{けい|さん}できないものをすべてえらび、その\R{理由}{り|ゆう}も\R{書}{か}きましょう。\par
ア\enspace$\qty{5}{m}+\qty{3}{m}$\hspace{2\zw}イ\enspace\qty{5}{m}＋3\R{秒}{びょう}\hspace{2\zw}ウ\enspace$\qty{2}{L}+\qty{500}{mL}$\hspace{2\zw}エ\enspace\qty{6}{kg}＋\qty{4}{\degreeCelsius}
\bkpans{イ、エ。\R{種類}{しゅ|るい}のちがう\R{量}{りょう}（\R{長}{なが}さと\R{時間}{じ|かん}、\R{重}{おも}さと\R{温度}{おん|ど}）はたせないから}

\bkproblem \R{表}{ひょう}は、ある\R{冬}{ふゆ}の\R{日}{ひ}の\R{気温}{き|おん}です。
\begin{bkcenter}
\begin{bktable}{colspec={Q[l,m,wd=24mm] *{5}{Q[c,m,wd=15mm]}}, rowsep=1.5pt}
\R{時刻}{じ|こく} & 6\R{時}{じ} & 9\R{時}{じ} & 12\R{時}{じ} & 15\R{時}{じ} & 18\R{時}{じ} \\
\R{気温}{き|おん}（℃） & $-3$ & 1 & 6 & 5 & 0 \\
\end{bktable}
\end{bkcenter}
（1）\nobreak 6\R{時}{じ}→12\R{時}{じ}、（2）\nobreak 12\R{時}{じ}→18\R{時}{じ}の\R{気温}{き|おん}の\R{変化量}{へん|か|りょう}を\R{求}{もと}めましょう。\par
（3）\nobreak \R{気温}{き|おん}がいちばん\R{大}{おお}きく\R{上}{あ}がったのは、どの3\R{時間}{じ|かん}でしょう。
\bkpans{（1）\nobreak $6-(-3)=9$\quad $+\qty{9}{\degreeCelsius}$\quad（2）\nobreak $0-6=-6$\quad $\qty{-6}{\degreeCelsius}$\quad（3）\nobreak 9\R{時}{じ}から12\R{時}{じ}（$+\qty{5}{\degreeCelsius}$）}

\bkproblem \R{電車}{でん|しゃ}が、\R{時刻}{じ|こく}10\R{秒}{びょう}にスタートから\qty{80}{m}、\R{時刻}{じ|こく}40\R{秒}{びょう}に\qty{530}{m}の\R{位置}{い|ち}にいました。$\Delta t$、$\Delta x$ と\R{速}{はや}さを\R{求}{もと}めましょう。
\begin{bkfigure}
\begin{tikzpicture}[x=0.2mm, y=1mm]
  \draw[bk axis] (-10,0) -- (625,0);
  \node[anchor=west, bk tick] at (628,0) {m};
  \foreach \p in {0,100,...,600} {\draw[bkInk, line width=0.6pt] (\p,-1.3) -- (\p,1.3);
    \node[bk tick, font=\sffamily\footnotesize, below=1.6mm] at (\p,0) {\p};}
  % 電車（先頭が位置の点）  % noruby
  \foreach \q/\c/\t in {80/bkGray/{\R{時刻}{じ|こく}10\R{秒}{びょう}}, 530/bkRoute2/{\R{時刻}{じ|こく}40\R{秒}{びょう}}} {
    \begin{scope}[shift={(\q,3)}, x=1mm, y=1mm]
      \fill[\c, rounded corners=1.2mm] (-15,0) rectangle (0,6.5);
      \foreach \w in {-13,-9,-5} {\fill[white] (\w,3) rectangle ++(2.6,2.3);}
      \fill[white] (-2.2,3) rectangle (-0.8,5.3);
    \end{scope}
    \fill[\c] (\q,0) circle[radius=1.3mm];
    \node[anchor=south, font=\sffamily\footnotesize, text=\c!80!black] at (\q-37.5,11.2) {\t};}
\end{tikzpicture}
\end{bkfigure}
\bkpans{$\Delta t=\qty{40}{s}-\qty{10}{s}=\qty{30}{s}$、$\Delta x=\qty{530}{m}-\qty{80}{m}=\qty{450}{m}$、$\qty{450}{m}\div\qty{30}{s}=\qty{15}{m/s}$}

\bkproblem \R{自転車}{じ|てん|しゃ}が\qty{6}{m/s}で25\R{秒}{びょう}\R{走}{はし}りました。\R{何}{なん}m\R{進}{すす}んだでしょう。
\bkpans{$6\times25=150$\quad\qty{150}{m}}

\bkproblem 1\R{時間}{じ|かん}あたり\qty{0.8}{kWh}の\R{電気}{でん|き}を\R{使}{つか}うエアコンを、6\R{時間}{じ|かん}つけました。\R{使}{つか}った\R{電気}{でん|き}の\R{量}{りょう}は\R{何}{なん}kWhでしょう。
\bkpans{$0.8\times6=4.8$\quad\qty{4.8}{kWh}}

\Needspace{52mm}\bkproblemgroup{チャレンジ}{2}

\bkproblem グラフは、あるラジコンカーの\R{速}{はや}さと\R{時間}{じ|かん}の\R{関係}{かん|けい}です。\par
\noindent\begin{minipage}[c]{0.52\linewidth}
（1）\nobreak はじめの3\R{秒}{びょう}で\R{進}{すす}んだ\R{道}{みち}のりは？\par
（2）\nobreak 5\R{秒}{びょう}\R{全部}{ぜん|ぶ}で\R{進}{すす}んだ\R{道}{みち}のりは？\par
（3）\nobreak （2）は、グラフのどこの\R{面積}{めん|せき}かな。ノートにグラフをかいて、\R{色}{いろ}をぬろう。
\end{minipage}\hfill
\begin{minipage}[c]{0.45\linewidth}\raggedleft
\begin{tikzpicture}
\begin{axis}[bkplot, width=48mm, height=34mm, xmin=0, xmax=5.9, ymin=0, ymax=5, xtick={1,2,3,4,5}, ytick={1,2,3,4},
  xlabel={\R{時間}{じ|かん}（s）}, ylabel={\R{速}{はや}さ（m/s）}]
  \addplot[bkcurve] coordinates {(0,1) (3,1)};
  \addplot[bkcurve] coordinates {(3,4) (5,4)};
  \draw[bkBlue, line width=0.8pt, densely dashed] (axis cs:3,1) -- (axis cs:3,4);
\end{axis}
\end{tikzpicture}
\end{minipage}
\bkpans{（1）\nobreak $1\times3=3$\quad\qty{3}{m}\quad（2）\nobreak $3+4\times2=3+8=11$\quad\qty{11}{m}\quad（3）\nobreak 0〜3\R{秒}{びょう}の\R{長方形}{ちょう|ほう|けい}（たて1、よこ3）と、3〜5\R{秒}{びょう}の\R{長方形}{ちょう|ほう|けい}（たて4、よこ2）を\R{合}{あ}わせた\R{部分}{ぶ|ぶん}}

\bkproblem \R{水}{みず}そうに\R{水}{みず}を\R{入}{い}れています。\R{時刻}{じ|こく}2\R{分}{ふん}に\qty{30}{L}、\R{時刻}{じ|こく}7\R{分}{ふん}に\qty{80}{L}でした。\R{時刻}{じ|こく}の\R{変化量}{へん|か|りょう}と\R{水}{みず}の\R{量}{りょう}の\R{変化量}{へん|か|りょう}を\R{求}{もと}め、1\R{分}{ぷん}あたり\R{何}{なん}Lふえたかを\R{求}{もと}めましょう。
\bkpans{\R{時刻}{じ|こく}の\R{変化量}{へん|か|りょう} $7-2=5$（\R{分}{ふん}）、\R{水}{みず}の\R{量}{りょう}の\R{変化量}{へん|か|りょう} $80-30=50$（L）、$50\div5=10$\quad 1\R{分}{ぷん}あたり\qty{10}{L}}

\bkproblem おふろの\R{水}{みず}をぬいています。\R{時刻}{じ|こく}0\R{分}{ふん}に\qty{180}{L}、\R{時刻}{じ|こく}6\R{分}{ぷん}に\qty{60}{L}でした。\R{水}{みず}の\R{量}{りょう}の\R{変化量}{へん|か|りょう}は\R{何}{なん}Lでしょう。また、1\R{分}{ぷん}あたり\R{何}{なん}Lへったでしょう。
\bkpans{$60-180=-120$\quad $\qty{-120}{L}$（\qty{120}{L}へった）。$120\div6=20$\quad 1\R{分}{ぷん}あたり\qty{20}{L}へった}

\begin{furikaeri}
\begin{checklist}
\item \R{同}{おな}じ\R{種類}{しゅ|るい}の\R{量}{りょう}を、\R{単位}{たん|い}をそろえてたしたりひいたりできる。\unskip\nobreak\hfill\penalty50\hskip1em\null\nobreak\hfill\mbox{\footnotesize\sffamily\color{bkTheme}→\pageref{bk:2-3:onaji}ページ}
\item ひき\R{算}{ざん}で\R{変化量}{へん|か|りょう}を\R{求}{もと}め、$\Delta$ を\R{使}{つか}って\R{書}{か}ける。\unskip\nobreak\hfill\penalty50\hskip1em\null\nobreak\hfill\mbox{\footnotesize\sffamily\color{bkTheme}→\pageref{bk:2-3:henka}ページ}
\item わり\R{算}{ざん}で\R{速}{はや}さを、かけ\R{算}{ざん}で\R{道}{みち}のりを\R{求}{もと}められる。\unskip\nobreak\hfill\penalty50\hskip1em\null\nobreak\hfill\mbox{\footnotesize\sffamily\color{bkTheme}→\pageref{bk:2-3:wari}ページ}
\end{checklist}
\end{furikaeri}

\bkfukidashi{sora}{つぎの2-4では、\R{分数}{ぶん|すう}と\R{小数}{しょう|すう}の\R{計算}{けい|さん}を\R{学}{まな}ぶよ。この\R{章}{しょう}の「1あたり」が、また\R{活躍}{かつ|やく}するよ。}
